% Test sample: a small Beamer talk with invented text (talk.pdf, notes in talk.pdfpc).
% It has what the tests need: slides with overlays (same page label), section start pages that Beamer
% labels like the slide before them (noframenumbering), notes only on the first overlay of a slide.
% Rebuild: pdflatex talk.tex (twice); talk.pdfpc is written by hand in PDFPC's JSON format.
\documentclass[aspectratio=169]{beamer}
\usetheme{default}
\setbeamertemplate{navigation symbols}{}
\AtBeginSection[]{\begin{frame}[noframenumbering]\centering\Large\insertsection\end{frame}}
\title{A vegetable garden on the balcony}
\author{Sample Speaker}
\date{}

\begin{document}

\begin{frame}
  \titlepage
\end{frame}

\begin{frame}{About this talk}
  Small pots, little sun and a few minutes a day are enough for a good harvest.
\end{frame}

\begin{frame}{Plan}
  \begin{itemize}
    \item Choosing the pots and the soil
    \item Sowing and watering
    \item Harvest and storage
  \end{itemize}
\end{frame}

\section{Pots and soil}

\begin{frame}{Pots}
  \begin{itemize}
    \item<1-> Terracotta keeps the roots cool in summer
    \item<2-> Plastic is light and holds the water longer
  \end{itemize}
\end{frame}

\begin{frame}{Soil}
  \begin{itemize}
    \item<1-> Universal soil mixed with some sand
    \item<2-> A layer of gravel at the bottom for drainage
    \item<3-> Compost once a month from spring to autumn
  \end{itemize}
\end{frame}

\begin{frame}{Light}
  Tomatoes and peppers want six hours of direct sun; lettuce and herbs grow in half shade.
\end{frame}

\section{Sowing and watering}

\begin{frame}{Sowing calendar}
  \begin{itemize}
    \item<1-> Basil and tomatoes in April, indoors
    \item<2-> Lettuce from March to September, every three weeks
  \end{itemize}
\end{frame}

\begin{frame}{Watering}
  Water early in the morning, at the foot of the plant, never on the leaves.
\end{frame}

\section{Harvest}

\begin{frame}{Harvest}
  Pick often: many plants produce more when the fruits are taken young.
\end{frame}

\begin{frame}{Questions}
  Thank you!
\end{frame}

\end{document}
